\documentclass[aspectratio=169,11pt]{beamer}

% ============================================================
% CLASE 10 - FEATURE ENGINEERING
% Diplomado de Python y Análisis de Datos
% Universidad Marista
%
% Compatible con Overleaf + pdfLaTeX
% ============================================================

% ============================================================
% CODIFICACIÓN
% ============================================================

\usepackage[utf8]{inputenc}
\usepackage[T1]{fontenc}
\usepackage[spanish,es-nodecimaldot]{babel}
\usepackage{lmodern}

% ============================================================
% PAQUETES
% ============================================================

\usepackage{amsmath}
\usepackage{amssymb}
\usepackage{booktabs}
\usepackage{array}
\usepackage{multicol}
\usepackage{graphicx}
\usepackage{xcolor}
\usepackage{tikz}
\usepackage{listings}

\usetikzlibrary{
    positioning,
    arrows.meta,
    shapes.geometric,
    calc
}

% ============================================================
% TEMA
% ============================================================

\usetheme{Madrid}
\usecolortheme{default}

\setbeamertemplate{navigation symbols}{}

% ============================================================
% COLORES
% ============================================================

\definecolor{azulmarista}{RGB}{18,67,115}
\definecolor{azulclaro}{RGB}{224,238,250}

\definecolor{verde}{RGB}{38,120,85}
\definecolor{verdeclaro}{RGB}{225,245,235}

\definecolor{naranja}{RGB}{225,125,35}
\definecolor{naranjaclaro}{RGB}{253,237,218}

\definecolor{morado}{RGB}{104,72,150}
\definecolor{moradoclaro}{RGB}{238,230,248}

\definecolor{gris}{RGB}{80,80,80}
\definecolor{grisclaro}{RGB}{245,245,245}

\definecolor{rojo}{RGB}{175,50,50}
\definecolor{rojoclaro}{RGB}{250,230,230}

% ============================================================
% COLORES BEAMER
% ============================================================

\setbeamercolor{structure}{
    fg=azulmarista
}

\setbeamercolor{frametitle}{
    fg=white,
    bg=azulmarista
}

\setbeamercolor{title}{
    fg=azulmarista
}

\setbeamercolor{subtitle}{
    fg=gris
}

\setbeamercolor{block title}{
    fg=white,
    bg=azulmarista
}

\setbeamercolor{block body}{
    bg=azulclaro
}

\setbeamercolor{block title example}{
    fg=white,
    bg=verde
}

\setbeamercolor{block body example}{
    bg=verdeclaro
}

\setbeamercolor{block title alerted}{
    fg=white,
    bg=rojo
}

\setbeamercolor{block body alerted}{
    bg=rojoclaro
}

% ============================================================
% PIE
% ============================================================

\setbeamertemplate{footline}{
    \leavevmode
    \hbox{
        \begin{beamercolorbox}[
            wd=.74\paperwidth,
            ht=2.7ex,
            dp=1.1ex,
            leftskip=1em
        ]{author in head/foot}
            Diplomado de Python y Análisis de Datos
        \end{beamercolorbox}

        \begin{beamercolorbox}[
            wd=.26\paperwidth,
            ht=2.7ex,
            dp=1.1ex,
            center
        ]{date in head/foot}
            \insertframenumber/\inserttotalframenumber
        \end{beamercolorbox}
    }
}

% ============================================================
% CÓDIGO
% ============================================================

\lstdefinestyle{pythonstyle}{
    language=Python,
    backgroundcolor=\color{grisclaro},
    basicstyle=\ttfamily\scriptsize,
    keywordstyle=\color{azulmarista}\bfseries,
    commentstyle=\color{verde},
    stringstyle=\color{morado},
    numberstyle=\tiny\color{gris},
    numbers=left,
    numbersep=5pt,
    breaklines=true,
    breakatwhitespace=false,
    keepspaces=true,
    showspaces=false,
    showstringspaces=false,
    showtabs=false,
    tabsize=4,
    frame=single,
    rulecolor=\color{azulmarista}
}

\lstset{
    style=pythonstyle
}

% ============================================================
% INFORMACIÓN
% ============================================================

\title{Feature Engineering}

\subtitle{
Clase 10\\
Transformar datos en mejores variables
}

\author{
Diplomado de Python y Análisis de Datos
}

\institute{
Universidad Marista
}

\date{2026}

% ============================================================
% DOCUMENTO
% ============================================================

\begin{document}

% ============================================================
% PORTADA
% ============================================================

\begin{frame}
    \titlepage
\end{frame}

% ============================================================
% OBJETIVOS
% ============================================================

\begin{frame}{Objetivos de la sesión}

Al finalizar la clase, el estudiante será capaz de:

\begin{itemize}

    \item comprender qué es Feature Engineering;

    \item identificar features útiles, inútiles y peligrosas;

    \item transformar variables numéricas;

    \item escalar variables;

    \item codificar variables categóricas;

    \item utilizar One-Hot Encoding;

    \item crear variables derivadas;

    \item construir ratios e interacciones;

    \item crear variables temporales;

    \item aplicar transformaciones logarítmicas;

    \item evitar Data Leakage durante preprocessing;

    \item utilizar ColumnTransformer;

    \item construir un Pipeline de Scikit-learn;

    \item comparar modelos antes y después del Feature Engineering.

\end{itemize}

\end{frame}

% ============================================================
% CRONOGRAMA
% ============================================================

\begin{frame}{Ruta de las 4 horas}

\scriptsize

\begin{table}
\centering

\begin{tabular}{p{2cm}p{9cm}}
\toprule
\textbf{Tiempo} & \textbf{Bloque} \\
\midrule

0:00--0:15 &
¿Qué es Feature Engineering? \\

0:15--0:35 &
Selección de variables y variables peligrosas \\

0:35--1:00 &
Escalamiento: StandardScaler y MinMaxScaler \\

1:00--1:25 &
Variables categóricas y One-Hot Encoding \\

1:25--1:45 &
Binning y transformaciones numéricas \\

1:45--1:55 &
Descanso \\

1:55--2:20 &
Interacciones, ratios y Polynomial Features \\

2:20--2:40 &
Variables temporales y de negocio \\

2:40--3:05 &
Data Leakage en preprocessing \\

3:05--3:25 &
ColumnTransformer + Pipeline \\

3:25--3:45 &
Live Coding: antes vs después \\

3:45--3:55 &
Feature Engineering Challenge \\

3:55--4:00 &
Conclusiones y puente a Regresión \\

\bottomrule
\end{tabular}

\end{table}

\end{frame}

% ============================================================
\section{Introducción}
% ============================================================

\begin{frame}{Hasta ahora...}

Ya sabemos:

\[
X
\rightarrow
Modelo
\rightarrow
\hat y
\]

Pero hay una pregunta que todavía no hemos estudiado profundamente.

\vspace{0.6cm}

\begin{alertblock}{Pregunta}
¿De dónde salen las variables que colocamos dentro de $X$?
\end{alertblock}

\end{frame}

% ------------------------------------------------------------

\begin{frame}{Un algoritmo solamente ve números}

Supongamos:

\begin{center}

\begin{tabular}{lll}
\toprule
Fecha & Zona & Precio \\
\midrule
2026-08-01 & Centro & 2.8M \\
2026-08-02 & Norte & 2.1M \\
\bottomrule
\end{tabular}

\end{center}

\vspace{0.5cm}

Para nosotros:

\begin{center}

``Centro''
\end{center}

tiene significado.

Para muchos algoritmos:

\begin{center}

\Large
necesitamos una representación numérica.
\end{center}

\end{frame}

% ------------------------------------------------------------

\begin{frame}{Feature Engineering}

Feature Engineering es el proceso de:

\begin{center}

\Large
crear, seleccionar o transformar variables
\end{center}

para representar mejor:

\begin{center}

\Large
la información relevante del problema.
\end{center}

\vspace{0.4cm}

Conceptualmente:

\[
\boxed{
Datos\ crudos
\rightarrow
Features
\rightarrow
Modelo
}
\]

\end{frame}

% ------------------------------------------------------------

\begin{frame}{Una buena representación importa}

Supongamos que tenemos:

\[
metros
\]

y:

\[
habitaciones
\]

Podemos crear:

\[
\boxed{
metros\_por\_habitacion
=
\frac{
metros
}{
habitaciones
}
}
\]

Esta nueva variable podría capturar:

\begin{center}

\Large
espacio disponible por habitación.
\end{center}

\end{frame}

% ------------------------------------------------------------

\begin{frame}{Feature Engineering no es inventar datos}

No estamos creando información de la nada.

Estamos:

\begin{itemize}

    \item reorganizando;

    \item transformando;

    \item representando;

    \item combinando;

    \item simplificando

\end{itemize}

información que ya existe.

\end{frame}

% ============================================================
\section{Selección de variables}
% ============================================================

\begin{frame}{¿Toda columna debe entrar al modelo?}

No.

Supongamos:

\begin{center}

\begin{tabular}{llll}
\toprule
id & nombre & metros & precio \\
\midrule
1001 & Ana & 120 & 2.5M \\
1002 & Pedro & 90 & 1.9M \\
\bottomrule
\end{tabular}

\end{center}

\vspace{0.5cm}

¿Debe entrar:

\[
id
\]

como feature?

No necesariamente.

\end{frame}

% ------------------------------------------------------------

\begin{frame}{Variables identificadoras}

Ejemplos:

\begin{itemize}

    \item ID de cliente;

    \item número de expediente;

    \item UUID;

    \item número de factura;

    \item nombre completo.

\end{itemize}

Normalmente sirven para:

\begin{center}

\Large
identificar
\end{center}

no necesariamente para:

\begin{center}

\Large
predecir.
\end{center}

\end{frame}

% ------------------------------------------------------------

\begin{frame}{Pregunta clave}

Antes de utilizar una columna:

\begin{center}

\Huge
¿Existe una razón por la cual esta variable
debería contener información predictiva?
\end{center}

\vspace{0.5cm}

No todas las columnas son features.

\end{frame}

% ============================================================
\section{Leakage}
% ============================================================

\begin{frame}{Variable peligrosa}

Queremos predecir:

\[
Churn
\]

y utilizamos:

\[
fecha\_cancelacion
\]

como feature.

\vspace{0.5cm}

\begin{alertblock}{Problema}
La variable prácticamente contiene la respuesta.
\end{alertblock}

Esto es:

\[
\boxed{
Data\ Leakage
}
\]

\end{frame}

% ------------------------------------------------------------

\begin{frame}{Regla de oro}

Para cada feature preguntar:

\begin{center}

\Large
¿Esta información existiría realmente
en el momento de hacer la predicción?
\end{center}

Si:

\[
No
\]

debemos investigar posible leakage.

\end{frame}

% ============================================================
\section{Variables numéricas}
% ============================================================

\begin{frame}{Escalas diferentes}

Supongamos:

\[
Edad=35
\]

pero:

\[
Ingreso=800\,000
\]

y:

\[
Visitas=5
\]

\vspace{0.5cm}

Las escalas son muy diferentes.

\end{frame}

% ------------------------------------------------------------

\begin{frame}{¿Por qué importa la escala?}

Algunos algoritmos utilizan:

\begin{itemize}

    \item distancia;

    \item gradientes;

    \item regularización.

\end{itemize}

Ejemplos sensibles:

\begin{itemize}

    \item KNN;

    \item K-Means;

    \item SVM;

    \item redes neuronales;

    \item modelos regularizados.

\end{itemize}

\end{frame}

% ------------------------------------------------------------

\begin{frame}{Ejemplo de distancia}

Tenemos dos diferencias:

\[
\Delta Edad=10
\]

\[
\Delta Ingreso=300\,000
\]

En distancia euclidiana:

\[
d
=
\sqrt{
10^2
+
300000^2
}
\]

\vspace{0.4cm}

¿Qué variable dominará?

\[
\boxed{Ingreso}
\]

\end{frame}

% ============================================================
\section{StandardScaler}
% ============================================================

\begin{frame}{Estandarización}

StandardScaler transforma aproximadamente:

\[
z
=
\frac{
x-\mu
}{
\sigma
}
\]

Después de transformar:

\begin{center}

media cercana a 0
\end{center}

y:

\begin{center}

desviación estándar cercana a 1.
\end{center}

\end{frame}

% ------------------------------------------------------------

\begin{frame}{Ejemplo}

Tenemos:

\[
x=120
\]

\[
\mu=100
\]

\[
\sigma=10
\]

Entonces:

\[
z
=
\frac{
120-100
}{
10
}
\]

\[
\boxed{
z=2
}
\]

\end{frame}

% ------------------------------------------------------------

\begin{frame}[fragile]{StandardScaler}

\begin{lstlisting}
from sklearn.preprocessing import (
    StandardScaler
)

scaler = StandardScaler()

X_train_scaled = scaler.fit_transform(
    X_train
)

X_test_scaled = scaler.transform(
    X_test
)
\end{lstlisting}

\end{frame}

% ------------------------------------------------------------

\begin{frame}{fit\_transform vs transform}

En Train:

\[
\boxed{
fit\_transform()
}
\]

porque:

\begin{center}

aprendemos media y desviación
\end{center}

y transformamos.

En Test:

\[
\boxed{
transform()
}
\]

porque utilizamos exactamente los parámetros aprendidos en Train.

\end{frame}

% ------------------------------------------------------------

\begin{frame}{Nunca hacemos esto}

\begin{center}

\Large
\texttt{scaler.fit\_transform(X\_test)}
\end{center}

\vspace{0.5cm}

¿Por qué?

Porque estaríamos permitiendo que Test influya en:

\[
\mu
\]

y:

\[
\sigma
\]

utilizados en su propia transformación.

\end{frame}

% ============================================================
\section{MinMaxScaler}
% ============================================================

\begin{frame}{Normalización Min-Max}

Otra transformación común:

\[
x'
=
\frac{
x-x_{\min}
}{
x_{\max}-x_{\min}
}
\]

Normalmente lleva los datos a:

\[
[0,1]
\]

\end{frame}

% ------------------------------------------------------------

\begin{frame}{Ejemplo Min-Max}

Supongamos:

\[
x_{\min}=0
\]

\[
x_{\max}=100
\]

y:

\[
x=40
\]

Entonces:

\[
x'
=
\frac{
40
}{
100
}
\]

\[
\boxed{
x'=0.4
}
\]

\end{frame}

% ------------------------------------------------------------

\begin{frame}[fragile]{MinMaxScaler}

\begin{lstlisting}
from sklearn.preprocessing import (
    MinMaxScaler
)

scaler = MinMaxScaler()

X_train_scaled = scaler.fit_transform(
    X_train
)

X_test_scaled = scaler.transform(
    X_test
)
\end{lstlisting}

\end{frame}

% ------------------------------------------------------------

\begin{frame}{StandardScaler vs MinMaxScaler}

\begin{center}

\begin{tabular}{lll}
\toprule
 & StandardScaler & MinMaxScaler \\
\midrule

Centro &
0 &
No necesariamente \\

Escala &
Desviación 1 &
$[0,1]$ \\

Outliers &
Puede afectarse &
Muy sensible \\

\bottomrule
\end{tabular}

\end{center}

\vspace{0.5cm}

No existe un escalador universalmente mejor.

\end{frame}

% ============================================================
\section{Árboles y escalamiento}
% ============================================================

\begin{frame}{¿Todos los modelos necesitan escalamiento?}

No.

Modelos basados en árboles como:

\begin{itemize}

    \item Decision Tree;

    \item Random Forest;

    \item Gradient Boosting

\end{itemize}

generalmente son mucho menos sensibles a la escala.

\vspace{0.4cm}

\begin{block}{Idea}
El preprocessing depende del algoritmo.
\end{block}

\end{frame}

% ============================================================
\section{Variables categóricas}
% ============================================================

\begin{frame}{El problema de las categorías}

Tenemos:

\[
Zona
\]

con:

\begin{center}

Centro

Norte

Sur
\end{center}

Muchos modelos requieren una representación numérica.

\end{frame}

% ------------------------------------------------------------

\begin{frame}{Una mala idea}

Podríamos escribir:

\[
Centro=1
\]

\[
Norte=2
\]

\[
Sur=3
\]

Pero esto podría sugerir:

\[
Sur>Norte>Centro
\]

aunque no exista ese orden.

\end{frame}

% ============================================================
\section{One-Hot Encoding}
% ============================================================

\begin{frame}{One-Hot Encoding}

Transformamos:

\[
Zona
\]

en variables binarias:

\begin{center}

\begin{tabular}{ccc}
\toprule
Centro & Norte & Sur \\
\midrule
1 & 0 & 0 \\
0 & 1 & 0 \\
0 & 0 & 1 \\
\bottomrule
\end{tabular}

\end{center}

\end{frame}

% ------------------------------------------------------------

\begin{frame}[fragile]{Con Pandas}

\begin{lstlisting}
df_encoded = pd.get_dummies(
    df,
    columns=[
        "zona"
    ],
    drop_first=True
)

print(
    df_encoded.head()
)
\end{lstlisting}

\end{frame}

% ------------------------------------------------------------

\begin{frame}[fragile]{Con Scikit-learn}

\begin{lstlisting}
from sklearn.preprocessing import (
    OneHotEncoder
)

encoder = OneHotEncoder(
    handle_unknown="ignore",
    sparse_output=False
)

encoder.fit(
    X_train[["zona"]]
)

zona_train = encoder.transform(
    X_train[["zona"]]
)

zona_test = encoder.transform(
    X_test[["zona"]]
)
\end{lstlisting}

\end{frame}

% ------------------------------------------------------------

\begin{frame}{handle\_unknown}

Supongamos que Train tiene:

\begin{center}

Centro, Norte, Sur
\end{center}

pero Test contiene:

\begin{center}

Poniente
\end{center}

\vspace{0.4cm}

Con:

\[
handle\_unknown="ignore"
\]

evitamos un error directo por esa categoría desconocida.

\end{frame}

% ============================================================
\section{Ordinal Encoding}
% ============================================================

\begin{frame}{Categorías con orden}

Algunas categorías sí tienen orden.

Ejemplo:

\begin{center}

Bajo

Medio

Alto
\end{center}

Aquí tiene sentido una representación como:

\[
Bajo=0
\]

\[
Medio=1
\]

\[
Alto=2
\]

\end{frame}

% ------------------------------------------------------------

\begin{frame}{Nominal vs ordinal}

\begin{columns}

\column{0.45\textwidth}

\begin{block}{Nominal}

Sin orden natural.

\begin{itemize}

    \item ciudad;

    \item color;

    \item producto.

\end{itemize}

\end{block}

\column{0.45\textwidth}

\begin{exampleblock}{Ordinal}

Con orden natural.

\begin{itemize}

    \item bajo/medio/alto;

    \item satisfacción;

    \item nivel educativo.

\end{itemize}

\end{exampleblock}

\end{columns}

\end{frame}

% ------------------------------------------------------------

\begin{frame}[fragile]{OrdinalEncoder}

\begin{lstlisting}
from sklearn.preprocessing import (
    OrdinalEncoder
)

encoder = OrdinalEncoder(
    categories=[
        [
            "Bajo",
            "Medio",
            "Alto"
        ]
    ]
)
\end{lstlisting}

\end{frame}

% ------------------------------------------------------------

\begin{frame}{¿Y LabelEncoder?}

\texttt{LabelEncoder} suele utilizarse principalmente para:

\[
y
\]

es decir:

\begin{center}

las etiquetas objetivo.
\end{center}

\vspace{0.4cm}

Para columnas categóricas dentro de $X$ normalmente preferimos:

\begin{itemize}

    \item OneHotEncoder;

    \item OrdinalEncoder.

\end{itemize}

\end{frame}

% ============================================================
\section{Binning}
% ============================================================

\begin{frame}{Binning}

Binning consiste en convertir una variable continua en intervalos.

Ejemplo:

\[
Edad
\]

se convierte en:

\begin{center}

Joven

Adulto

Mayor
\end{center}

\end{frame}

% ------------------------------------------------------------

\begin{frame}[fragile]{Binning manual}

\begin{lstlisting}
df["grupo_edad"] = pd.cut(
    df["edad"],
    bins=[
        0,
        25,
        40,
        60,
        100
    ],
    labels=[
        "18-25",
        "26-40",
        "41-60",
        "60+"
    ]
)
\end{lstlisting}

\end{frame}

% ------------------------------------------------------------

\begin{frame}{¿Por qué utilizar binning?}

Puede ayudar a:

\begin{itemize}

    \item simplificar relaciones;

    \item crear segmentos interpretables;

    \item capturar umbrales;

    \item incorporar conocimiento del negocio.

\end{itemize}

Pero:

\begin{alertblock}{Cuidado}
También perdemos información al convertir números continuos
en categorías.
\end{alertblock}

\end{frame}

% ============================================================
\section{Transformación logarítmica}
% ============================================================

\begin{frame}{Variables muy sesgadas}

Supongamos ingresos:

\[
10000,\,
15000,\,
20000,\,
25000,\,
1000000
\]

Existe una cola muy larga.

Podemos estudiar:

\[
\log(1+x)
\]

\end{frame}

% ------------------------------------------------------------

\begin{frame}[fragile]{Transformación log}

\begin{lstlisting}
import numpy as np

df["ingreso_log"] = np.log1p(
    df["ingreso"]
)
\end{lstlisting}

\texttt{log1p(x)} calcula:

\[
\log(1+x)
\]

y permite manejar:

\[
x=0
\]

sin obtener $\log(0)$.

\end{frame}

% ------------------------------------------------------------

\begin{frame}{¿Qué logra el log?}

Puede:

\begin{itemize}

    \item reducir asimetría;

    \item comprimir valores extremos;

    \item facilitar algunas relaciones;

    \item estabilizar escalas.

\end{itemize}

\vspace{0.4cm}

\begin{alertblock}{Cuidado}
No debemos transformar variables automáticamente sin una razón.
\end{alertblock}

\end{frame}

% ============================================================
% DESCANSO
% ============================================================

\begin{frame}

\begin{center}

{\Huge
\textcolor{azulmarista}{
\textbf{DESCANSO}
}
}

\vspace{0.6cm}

{\Large
10 minutos
}

\vspace{1cm}

{\Large
Al regresar:
}

\vspace{0.4cm}

{\LARGE
¿Cómo crear información nueva a partir
de las variables existentes?
}

\end{center}

\end{frame}

% ============================================================
\section{Interacciones}
% ============================================================

\begin{frame}{Una variable puede depender de otra}

Supongamos:

\[
metros
\]

y:

\[
zona
\]

Quizá:

\[
10\text{ m}^2
\]

adicionales no valen lo mismo:

\begin{center}

en Centro
\end{center}

que:

\begin{center}

en una zona periférica.
\end{center}

\end{frame}

% ------------------------------------------------------------

\begin{frame}{Interacción}

Una interacción estudia cómo el efecto de una variable puede depender
de otra.

Ejemplo:

\[
x_1x_2
\]

Podemos crear:

\[
metros
\times
habitaciones
\]

aunque la interpretación depende del problema.

\end{frame}

% ============================================================
\section{Ratios}
% ============================================================

\begin{frame}{Ratios}

Muchas features útiles provienen de razones.

Ejemplos:

\[
metros\_por\_habitacion
=
\frac{
metros
}{
habitaciones
}
\]

\[
gasto\_promedio
=
\frac{
gasto\_total
}{
numero\_compras
}
\]

\end{frame}

% ------------------------------------------------------------

\begin{frame}[fragile]{Crear ratios}

\begin{lstlisting}
df["metros_por_habitacion"] = (
    df["metros"]
    /
    df["habitaciones"]
)

df["gasto_promedio"] = (
    df["gasto_total"]
    /
    df["numero_compras"]
)
\end{lstlisting}

\end{frame}

% ------------------------------------------------------------

\begin{frame}{Cuidado con división entre cero}

Si:

\[
numero\_compras=0
\]

entonces:

\[
\frac{
gasto
}{
numero\_compras
}
\]

no está definido.

\vspace{0.4cm}

Siempre debemos validar las reglas de negocio.

\end{frame}

% ============================================================
\section{Polynomial Features}
% ============================================================

\begin{frame}{Relaciones no lineales}

Supongamos que:

\[
y
\]

no cambia linealmente con:

\[
x
\]

Podríamos incluir:

\[
x^2
\]

o:

\[
x_1x_2
\]

como features adicionales.

\end{frame}

% ------------------------------------------------------------

\begin{frame}{Ejemplo polinomial}

Un modelo:

\[
y
=
\beta_0
+
\beta_1x
+
\beta_2x^2
\]

puede representar curvatura.

\vspace{0.5cm}

Seguimos utilizando un modelo lineal en los parámetros,
pero con variables transformadas.

\end{frame}

% ------------------------------------------------------------

\begin{frame}[fragile]{PolynomialFeatures}

\begin{lstlisting}
from sklearn.preprocessing import (
    PolynomialFeatures
)

poly = PolynomialFeatures(
    degree=2,
    include_bias=False
)

X_poly = poly.fit_transform(
    X[
        [
            "metros",
            "antiguedad"
        ]
    ]
)
\end{lstlisting}

\end{frame}

% ------------------------------------------------------------

\begin{frame}{¿Qué crea degree=2?}

Para:

\[
x_1,x_2
\]

puede crear:

\[
x_1
\]

\[
x_2
\]

\[
x_1^2
\]

\[
x_1x_2
\]

\[
x_2^2
\]

\vspace{0.4cm}

Esto aumenta la capacidad del modelo.

\end{frame}

% ------------------------------------------------------------

\begin{frame}{Más features también pueden ser peligrosas}

Si agregamos demasiadas variables:

\[
p\uparrow
\]

podemos aumentar:

\begin{itemize}

    \item complejidad;

    \item ruido;

    \item costo;

    \item riesgo de overfitting.

\end{itemize}

\vspace{0.4cm}

\begin{block}{Regla}
Más features no significa automáticamente mejor modelo.
\end{block}

\end{frame}

% ============================================================
\section{Variables temporales}
% ============================================================

\begin{frame}{Una fecha contiene muchas features}

Tenemos:

\[
2026-08-21
\]

Podemos obtener:

\begin{itemize}

    \item año;

    \item mes;

    \item día;

    \item día de semana;

    \item trimestre;

    \item fin de semana.

\end{itemize}

\end{frame}

% ------------------------------------------------------------

\begin{frame}[fragile]{Features temporales}

\begin{lstlisting}
df["fecha"] = pd.to_datetime(
    df["fecha"]
)

df["anio"] = (
    df["fecha"].dt.year
)

df["mes"] = (
    df["fecha"].dt.month
)

df["dia_semana"] = (
    df["fecha"].dt.dayofweek
)

df["trimestre"] = (
    df["fecha"].dt.quarter
)
\end{lstlisting}

\end{frame}

% ------------------------------------------------------------

\begin{frame}[fragile]{Fin de semana}

\begin{lstlisting}
df["fin_semana"] = (
    df["fecha"]
    .dt.dayofweek
    .isin([
        5,
        6
    ])
    .astype(int)
)
\end{lstlisting}

Ahora:

\[
1
\]

significa fin de semana.

\end{frame}

% ------------------------------------------------------------

\begin{frame}{Conocimiento del negocio}

Feature Engineering también utiliza conocimiento del dominio.

Ejemplos:

\begin{itemize}

    \item antigüedad del cliente;

    \item días desde última compra;

    \item gasto promedio;

    \item tasa de conversión;

    \item frecuencia de visitas;

    \item ocupación por habitación.

\end{itemize}

\end{frame}

% ------------------------------------------------------------

\begin{frame}{RFM como ejemplo}

En clientes podemos construir:

\[
R
=
Recency
\]

\[
F
=
Frequency
\]

\[
M
=
Monetary
\]

Estas tres variables resumen:

\begin{itemize}

    \item qué tan reciente compró;

    \item qué tan seguido compra;

    \item cuánto gasta.

\end{itemize}

\end{frame}

% ============================================================
\section{Preprocessing Leakage}
% ============================================================

\begin{frame}{Leakage también ocurre al transformar}

Supongamos que calculamos:

\[
\mu
\]

y:

\[
\sigma
\]

utilizando:

\[
Train+Test
\]

antes de escalar.

\vspace{0.5cm}

Entonces Test influyó en el preprocessing.

\[
\boxed{
Leakage
}
\]

\end{frame}

% ------------------------------------------------------------

\begin{frame}{Orden correcto}

Primero:

\[
Train/Test
\]

Después:

\[
fit\ preprocessing\ en\ Train
\]

Después:

\[
transform\ Train
\]

y:

\[
transform\ Test
\]

utilizando solamente lo aprendido en Train.

\end{frame}

% ------------------------------------------------------------

\begin{frame}{Flujo incorrecto}

\[
Datos
\]

\[
\downarrow
\]

\[
Scaler.fit(Datos)
\]

\[
\downarrow
\]

\[
Train/Test
\]

\vspace{0.4cm}

\begin{alertblock}{Problema}
El scaler ya vio información que terminará en Test.
\end{alertblock}

\end{frame}

% ------------------------------------------------------------

\begin{frame}{Flujo correcto}

\[
Datos
\]

\[
\downarrow
\]

\[
Train/Test
\]

\[
\downarrow
\]

\[
Scaler.fit(Train)
\]

\[
\downarrow
\]

\[
Transform(Train)
\qquad
Transform(Test)
\]

\end{frame}

% ============================================================
\section{ColumnTransformer}
% ============================================================

\begin{frame}{Datasets mixtos}

Nuestro dataset puede contener:

\begin{center}

\begin{tabular}{ll}
\toprule
Variable & Tipo \\
\midrule
metros & numérica \\
edad & numérica \\
zona & categórica \\
tipo & categórica \\
\bottomrule
\end{tabular}

\end{center}

Necesitamos diferentes transformaciones.

\end{frame}

% ------------------------------------------------------------

\begin{frame}{La solución}

Podemos decir:

\begin{center}

Variables numéricas
\end{center}

\[
\rightarrow
StandardScaler
\]

y:

\begin{center}

Variables categóricas
\end{center}

\[
\rightarrow
OneHotEncoder
\]

utilizando:

\[
\boxed{
ColumnTransformer
}
\]

\end{frame}

% ------------------------------------------------------------

\begin{frame}[fragile]{Definir columnas}

\begin{lstlisting}
numericas = [
    "metros",
    "habitaciones",
    "banios",
    "antiguedad"
]

categoricas = [
    "zona",
    "tipo_vivienda"
]
\end{lstlisting}

\end{frame}

% ------------------------------------------------------------

\begin{frame}[fragile]{ColumnTransformer}

\begin{lstlisting}
from sklearn.compose import (
    ColumnTransformer
)

from sklearn.preprocessing import (
    StandardScaler,
    OneHotEncoder
)

preprocesador = ColumnTransformer(
    transformers=[
        (
            "num",
            StandardScaler(),
            numericas
        ),
        (
            "cat",
            OneHotEncoder(
                handle_unknown="ignore"
            ),
            categoricas
        )
    ]
)
\end{lstlisting}

\end{frame}

% ============================================================
\section{Pipeline}
% ============================================================

\begin{frame}{Pipeline}

Un Pipeline permite conectar:

\[
Preprocessing
\]

con:

\[
Modelo
\]

para evitar procesos manuales inconsistentes.

Conceptualmente:

\[
X
\rightarrow
Transformar
\rightarrow
Modelo
\rightarrow
Predicción
\]

\end{frame}

% ------------------------------------------------------------

\begin{frame}[fragile]{Pipeline completo}

\begin{lstlisting}
from sklearn.pipeline import Pipeline
from sklearn.linear_model import (
    LinearRegression
)

pipeline = Pipeline(
    steps=[
        (
            "preprocesamiento",
            preprocesador
        ),
        (
            "modelo",
            LinearRegression()
        )
    ]
)
\end{lstlisting}

\end{frame}

% ------------------------------------------------------------

\begin{frame}[fragile]{Entrenar Pipeline}

\begin{lstlisting}
pipeline.fit(
    X_train,
    y_train
)

predicciones = pipeline.predict(
    X_test
)
\end{lstlisting}

\vspace{0.4cm}

El Pipeline se encarga de:

\begin{enumerate}

    \item aprender preprocessing en Train;

    \item transformar;

    \item entrenar el modelo;

    \item transformar correctamente Test;

    \item predecir.

\end{enumerate}

\end{frame}

% ------------------------------------------------------------

\begin{frame}{Ventaja fundamental}

Con Pipeline reducimos el riesgo de:

\begin{itemize}

    \item olvidar transformaciones;

    \item transformar Train y Test diferente;

    \item introducir leakage accidental;

    \item romper el flujo al llevar el modelo a producción.

\end{itemize}

\end{frame}

% ============================================================
\section{Live Coding}
% ============================================================

\begin{frame}{Caso práctico}

Utilizaremos:

\begin{center}

\Large
\texttt{viviendas\_features.csv}
\end{center}

Variables:

\begin{itemize}

    \item metros;

    \item habitaciones;

    \item baños;

    \item antigüedad;

    \item zona;

    \item tipo\_vivienda;

    \item fecha\_publicacion;

    \item precio.

\end{itemize}

\end{frame}

% ------------------------------------------------------------

\begin{frame}{Objetivo}

Compararemos dos experimentos.

\[
\boxed{
Modelo\ A
}
\]

con variables básicas.

Contra:

\[
\boxed{
Modelo\ B
}
\]

con Feature Engineering.

\end{frame}

% ------------------------------------------------------------

\begin{frame}[fragile]{Paso 1: cargar}

\begin{lstlisting}
import pandas as pd

df = pd.read_csv(
    "viviendas_features.csv"
)

print(
    df.head()
)

df.info()
\end{lstlisting}

\end{frame}

% ------------------------------------------------------------

\begin{frame}[fragile]{Paso 2: limpiar fecha}

\begin{lstlisting}
df["fecha_publicacion"] = (
    pd.to_datetime(
        df["fecha_publicacion"]
    )
)
\end{lstlisting}

\end{frame}

% ------------------------------------------------------------

\begin{frame}[fragile]{Paso 3: variables temporales}

\begin{lstlisting}
df["mes_publicacion"] = (
    df["fecha_publicacion"]
    .dt.month
)

df["dia_semana"] = (
    df["fecha_publicacion"]
    .dt.dayofweek
)
\end{lstlisting}

\end{frame}

% ------------------------------------------------------------

\begin{frame}[fragile]{Paso 4: nuevas features}

\begin{lstlisting}
df["metros_por_habitacion"] = (
    df["metros"]
    /
    df["habitaciones"]
)

df["banios_por_habitacion"] = (
    df["banios"]
    /
    df["habitaciones"]
)
\end{lstlisting}

\end{frame}

% ------------------------------------------------------------

\begin{frame}{Validar antes de dividir}

Antes de crear ratios revisar:

\[
habitaciones=0?
\]

También:

\[
banios<0?
\]

\[
metros\leq0?
\]

\vspace{0.4cm}

Feature Engineering no reemplaza la validación de datos.

\end{frame}

% ------------------------------------------------------------

\begin{frame}[fragile]{Paso 5: X e y}

\begin{lstlisting}
X = df.drop(
    columns=[
        "precio",
        "fecha_publicacion"
    ]
)

y = df[
    "precio"
]
\end{lstlisting}

\end{frame}

% ------------------------------------------------------------

\begin{frame}[fragile]{Paso 6: Train/Test}

\begin{lstlisting}
from sklearn.model_selection import (
    train_test_split
)

X_train, X_test, y_train, y_test = (
    train_test_split(
        X,
        y,
        test_size=0.20,
        random_state=42
    )
)
\end{lstlisting}

\end{frame}

% ------------------------------------------------------------

\begin{frame}[fragile]{Paso 7: columnas}

\begin{lstlisting}
numericas = [
    "metros",
    "habitaciones",
    "banios",
    "antiguedad",
    "mes_publicacion",
    "dia_semana",
    "metros_por_habitacion",
    "banios_por_habitacion"
]

categoricas = [
    "zona",
    "tipo_vivienda"
]
\end{lstlisting}

\end{frame}

% ------------------------------------------------------------

\begin{frame}[fragile]{Paso 8: preprocessing}

\begin{lstlisting}
preprocesador = ColumnTransformer(
    transformers=[
        (
            "numericas",
            StandardScaler(),
            numericas
        ),
        (
            "categoricas",
            OneHotEncoder(
                handle_unknown="ignore"
            ),
            categoricas
        )
    ]
)
\end{lstlisting}

\end{frame}

% ------------------------------------------------------------

\begin{frame}[fragile]{Paso 9: pipeline}

\begin{lstlisting}
pipeline = Pipeline(
    steps=[
        (
            "prep",
            preprocesador
        ),
        (
            "modelo",
            LinearRegression()
        )
    ]
)

pipeline.fit(
    X_train,
    y_train
)
\end{lstlisting}

\end{frame}

% ------------------------------------------------------------

\begin{frame}[fragile]{Paso 10: predecir}

\begin{lstlisting}
pred = pipeline.predict(
    X_test
)
\end{lstlisting}

Ahora:

\[
\hat y
\]

fue producido después de aplicar automáticamente todas
las transformaciones correspondientes.

\end{frame}

% ------------------------------------------------------------

\begin{frame}[fragile]{Paso 11: evaluar}

\begin{lstlisting}
from sklearn.metrics import (
    mean_absolute_error,
    mean_squared_error,
    r2_score
)

mae = mean_absolute_error(
    y_test,
    pred
)

rmse = (
    mean_squared_error(
        y_test,
        pred
    ) ** 0.5
)

r2 = r2_score(
    y_test,
    pred
)

print(
    mae,
    rmse,
    r2
)
\end{lstlisting}

\end{frame}

% ============================================================
\section{Antes vs después}
% ============================================================

\begin{frame}{Experimento}

Modelo inicial:

\[
X=
\{
metros,
habitaciones,
banios,
antiguedad
\}
\]

Modelo mejorado:

\[
X=
\{
variables\ originales
+
categorias
+
ratios
+
fecha
\}
\]

\end{frame}

% ------------------------------------------------------------

\begin{frame}{Comparar}

Ejemplo:

\begin{center}

\begin{tabular}{lrr}
\toprule
Modelo & MAE & $R^2$ \\
\midrule

Original & 220000 & 0.72 \\

Feature Engineering & 145000 & 0.86 \\

\bottomrule
\end{tabular}

\end{center}

\vspace{0.5cm}

\begin{exampleblock}{Pregunta}
¿Qué variables nuevas pudieron aportar información?
\end{exampleblock}

\end{frame}

% ------------------------------------------------------------

\begin{frame}{Pero también puede empeorar}

Supongamos:

\[
MAE_{antes}=150000
\]

\[
MAE_{despues}=175000
\]

\vspace{0.5cm}

\begin{alertblock}{Importante}
Una feature nueva debe demostrar que aporta valor.
\end{alertblock}

No debemos conservar variables solamente porque parecen interesantes.

\end{frame}

% ============================================================
\section{Feature Importance}
% ============================================================

\begin{frame}{¿Qué variables utiliza más un modelo?}

Algunos modelos ofrecen medidas de:

\[
\boxed{
Feature\ Importance
}
\]

Por ejemplo:

\begin{itemize}

    \item árboles;

    \item Random Forest;

    \item Gradient Boosting.

\end{itemize}

\end{frame}

% ------------------------------------------------------------

\begin{frame}[fragile]{Ejemplo rápido}

\begin{lstlisting}
from sklearn.ensemble import (
    RandomForestRegressor
)

modelo_rf = RandomForestRegressor(
    n_estimators=100,
    random_state=42
)

modelo_rf.fit(
    X_train_num,
    y_train
)

print(
    modelo_rf.feature_importances_
)
\end{lstlisting}

\end{frame}

% ------------------------------------------------------------

\begin{frame}{Cuidado con importancia}

Importancia no significa:

\[
\boxed{
causalidad
}
\]

Además, diferentes métodos de importancia tienen diferentes propiedades.

\vspace{0.4cm}

Más adelante veremos:

\[
\boxed{
SHAP
}
\]

para explicabilidad.

\end{frame}

% ============================================================
\section{Feature Engineering Challenge}
% ============================================================

\begin{frame}{Reto principal}

\begin{center}

{\Huge
\textcolor{azulmarista}{
\textbf{Feature Engineering Challenge}
}
}

\vspace{0.7cm}

{\Large
¿Podemos mejorar un modelo
sin cambiar el algoritmo?
}

\end{center}

\end{frame}

% ------------------------------------------------------------

\begin{frame}{Regla del reto}

Todos utilizarán:

\[
\boxed{
LinearRegression
}
\]

No pueden cambiar el algoritmo.

Solamente pueden modificar:

\[
\boxed{
X
}
\]

\vspace{0.4cm}

El objetivo:

\[
MAE\downarrow
\]

\end{frame}

% ------------------------------------------------------------

\begin{frame}{Dataset}

\begin{center}

\Large
\texttt{clientes\_features.csv}
\end{center}

Variables:

\begin{itemize}

    \item edad;

    \item ingreso;

    \item visitas;

    \item compras;

    \item gasto\_total;

    \item ciudad;

    \item fecha\_registro;

    \item gasto\_siguiente\_mes.

\end{itemize}

\end{frame}

% ------------------------------------------------------------

\begin{frame}{Target}

Queremos predecir:

\[
\boxed{
gasto\_siguiente\_mes
}
\]

Por lo tanto:

\[
y
=
gasto\_siguiente\_mes
\]

\end{frame}

% ------------------------------------------------------------

\begin{frame}{Fase 1}

Construir modelo base usando solamente:

\begin{itemize}

    \item edad;

    \item ingreso;

    \item visitas;

    \item compras;

    \item gasto\_total.

\end{itemize}

Registrar:

\[
MAE_{base}
\]

\end{frame}

% ------------------------------------------------------------

\begin{frame}{Fase 2}

Crear al menos:

\begin{enumerate}

    \item una variable ratio;

    \item una variable temporal;

    \item una transformación;

    \item una variable categórica codificada.

\end{enumerate}

\end{frame}

% ------------------------------------------------------------

\begin{frame}{Ideas posibles}

\[
gasto\_por\_compra
=
\frac{
gasto\_total
}{
compras
}
\]

\[
visitas\_por\_compra
=
\frac{
visitas
}{
compras
}
\]

\[
antiguedad\_cliente
\]

\[
ingreso\_log
\]

\[
ciudad\rightarrow OneHot
\]

\end{frame}

% ------------------------------------------------------------

\begin{frame}{Fase 3}

Crear:

\[
Pipeline
\]

con:

\begin{itemize}

    \item StandardScaler para numéricas;

    \item OneHotEncoder para categóricas;

    \item LinearRegression.

\end{itemize}

\end{frame}

% ------------------------------------------------------------

\begin{frame}{Fase 4}

Comparar:

\[
MAE_{base}
\]

contra:

\[
MAE_{features}
\]

Calcular:

\[
Mejora
=
MAE_{base}
-
MAE_{features}
\]

\end{frame}

% ------------------------------------------------------------

\begin{frame}{Fase 5}

Cada equipo deberá responder:

\begin{enumerate}

    \item ¿Qué feature nueva ayudó más?

    \item ¿Qué feature no ayudó?

    \item ¿Alguna produjo overfitting?

    \item ¿Existe riesgo de leakage?

    \item ¿Qué variables conservarían?

\end{enumerate}

\end{frame}

% ============================================================
\section{Errores comunes}
% ============================================================

\begin{frame}{Error 1}

\begin{center}

\Large
``Voy a escalar todo el dataset
antes del Train/Test.''
\end{center}

\vspace{0.5cm}

\begin{alertblock}{Incorrecto}
Puede introducir Data Leakage.
\end{alertblock}

\end{frame}

% ------------------------------------------------------------

\begin{frame}{Error 2}

\begin{center}

\Large
``Centro=1, Norte=2, Sur=3.''
\end{center}

\vspace{0.5cm}

\begin{alertblock}{Problema}
Estamos introduciendo un orden artificial
si la variable es nominal.
\end{alertblock}

\end{frame}

% ------------------------------------------------------------

\begin{frame}{Error 3}

\begin{center}

\Large
``Cuantas más features tenga,
mejor será el modelo.''
\end{center}

\vspace{0.5cm}

\begin{alertblock}{Incorrecto}
Las variables irrelevantes pueden aumentar ruido y complejidad.
\end{alertblock}

\end{frame}

% ------------------------------------------------------------

\begin{frame}{Error 4}

\begin{center}

\Large
``Si una variable mejora mucho el modelo,
debe ser excelente.''
\end{center}

\vspace{0.5cm}

Primero preguntar:

\begin{center}

¿Existe en el momento real de predicción?
\end{center}

Puede ser:

\[
\boxed{
Leakage
}
\]

\end{frame}

% ------------------------------------------------------------

\begin{frame}{Error 5}

\begin{center}

\Large
``Todos los algoritmos necesitan StandardScaler.''
\end{center}

\vspace{0.5cm}

No.

El preprocessing depende de:

\begin{itemize}

    \item algoritmo;

    \item distribución;

    \item objetivo;

    \item tipo de datos.

\end{itemize}

\end{frame}

% ============================================================
\section{Pipeline conceptual}
% ============================================================

\begin{frame}{Pipeline de Feature Engineering}

\begin{center}

\begin{tikzpicture}[
node distance=0.28cm,
every node/.style={
draw,
rounded corners,
fill=azulclaro,
minimum width=3cm,
minimum height=0.43cm,
font=\scriptsize
}
]

\node (raw) {Datos crudos};

\node[below=of raw] (split) {Train/Test};

\node[below=of split] (num) {Transformar numéricas};

\node[below=of num] (cat) {Codificar categóricas};

\node[below=of cat] (create) {Crear features};

\node[below=of create] (pipe) {Pipeline};

\node[below=of pipe] (model) {Modelo};

\node[below=of model] (metric) {Evaluar};

\node[below=of metric] (compare) {Comparar};

\node[below=of compare] (keep) {Conservar mejoras};

\draw[-{Latex}] (raw) -- (split);
\draw[-{Latex}] (split) -- (num);
\draw[-{Latex}] (num) -- (cat);
\draw[-{Latex}] (cat) -- (create);
\draw[-{Latex}] (create) -- (pipe);
\draw[-{Latex}] (pipe) -- (model);
\draw[-{Latex}] (model) -- (metric);
\draw[-{Latex}] (metric) -- (compare);
\draw[-{Latex}] (compare) -- (keep);

\end{tikzpicture}

\end{center}

\end{frame}

% ============================================================
\section{Herramientas}
% ============================================================

\begin{frame}{Herramientas que vimos}

\scriptsize

\begin{table}
\centering

\begin{tabular}{ll}
\toprule
\textbf{Objetivo} & \textbf{Herramienta} \\
\midrule

Estandarizar &
\texttt{StandardScaler()} \\

Normalizar &
\texttt{MinMaxScaler()} \\

One-Hot &
\texttt{OneHotEncoder()} \\

Ordinal &
\texttt{OrdinalEncoder()} \\

Binning &
\texttt{pd.cut()} \\

Log &
\texttt{np.log1p()} \\

Polinomios &
\texttt{PolynomialFeatures()} \\

Transformar columnas &
\texttt{ColumnTransformer()} \\

Pipeline &
\texttt{Pipeline()} \\

\bottomrule
\end{tabular}

\end{table}

\end{frame}

% ============================================================
\section{Ideas clave}
% ============================================================

\begin{frame}{15 ideas que deben recordar}

\scriptsize

\begin{enumerate}

    \item Los modelos aprenden de la representación que les damos.

    \item No toda columna debe ser una feature.

    \item IDs normalmente no aportan significado predictivo directo.

    \item Siempre debemos investigar leakage.

    \item La escala importa para algunos algoritmos.

    \item StandardScaler centra y escala.

    \item MinMaxScaler lleva normalmente a $[0,1]$.

    \item One-Hot Encoding sirve para categorías nominales.

    \item Ordinal Encoding requiere un orden real.

    \item Binning simplifica pero pierde información.

    \item Transformaciones logarítmicas pueden reducir asimetría.

    \item Ratios pueden capturar información de negocio.

    \item Interacciones permiten representar relaciones más complejas.

    \item El preprocessing debe aprenderse solamente con Train.

    \item Pipeline ayuda a construir procesos reproducibles y seguros.

\end{enumerate}

\end{frame}

% ============================================================
\section{Pregunta final}
% ============================================================

\begin{frame}{Pregunta de salida 1}

Tenemos:

\[
Edad=30
\]

\[
Ingreso=800000
\]

y utilizaremos:

\[
KNN
\]

\vspace{0.5cm}

\begin{alertblock}{Pregunta}
¿Qué transformación investigaríamos primero?
\end{alertblock}

Una posibilidad:

\[
\boxed{
Escalamiento
}
\]

\end{frame}

% ------------------------------------------------------------

\begin{frame}{Pregunta de salida 2}

Tenemos:

\begin{center}

Centro

Norte

Sur
\end{center}

sin orden natural.

\vspace{0.4cm}

¿Qué utilizaríamos?

\[
\boxed{
One-Hot\ Encoding
}
\]

\end{frame}

% ------------------------------------------------------------

\begin{frame}{Pregunta de salida 3}

Calculamos:

\[
StandardScaler.fit(X)
\]

antes de separar Train/Test.

¿Qué problema existe?

\[
\boxed{
Data\ Leakage
}
\]

\end{frame}

% ============================================================
\section{Puente a Regresión}
% ============================================================

\begin{frame}{Ya tenemos mejores variables}

Ahora podemos construir:

\[
X
\]

de mejor calidad.

Pero todavía falta comprender profundamente:

\begin{center}

\Huge
¿cómo aprende una regresión?
\end{center}

\end{frame}

% ------------------------------------------------------------

\begin{frame}{Próxima clase}

\begin{center}

{\Huge
\textbf{Clase 11}
}

\vspace{0.4cm}

{\LARGE
\textcolor{azulmarista}{
\textbf{Regresión}
}
}

\vspace{0.3cm}

{\Large
Predicción de variables numéricas
}

\end{center}

\vspace{0.5cm}

Veremos:

\begin{itemize}

    \item regresión lineal simple;

    \item mínimos cuadrados;

    \item pendiente e intercepto;

    \item regresión múltiple;

    \item interpretación de coeficientes;

    \item supuestos;

    \item residuos;

    \item MAE;

    \item RMSE;

    \item $R^2$;

    \item regularización;

    \item Ridge;

    \item Lasso.

\end{itemize}

\end{frame}

% ============================================================
% CIERRE
% ============================================================

\begin{frame}

\begin{center}

{\Huge
\textcolor{azulmarista}{
\textbf{Feature Engineering}
}
}

\vspace{0.8cm}

{\Large
El algoritmo importa.
}

\vspace{0.4cm}

{\Large
Pero también importa
}

\vspace{0.3cm}

{\LARGE
\textbf{cómo representamos el problema.}
}

\vspace{0.8cm}

\textbf{Diplomado de Python y Análisis de Datos}

\vspace{0.2cm}

\textbf{Universidad Marista}

\end{center}

\end{frame}

% ============================================================

\end{document}